\documentclass[11pt,a4paper]{article}

\usepackage[utf8]{inputenc}
\usepackage[spanish]{babel}
\usepackage[T1]{fontenc}
\usepackage[margin=2.7cm]{geometry}
\usepackage{amsmath,amssymb,mathtools}
\usepackage{graphicx}
\usepackage{booktabs}
\usepackage{algorithm}
\usepackage{algpseudocode}
\usepackage{fancyhdr}
\usepackage{xcolor}
\usepackage[colorlinks=true,linkcolor=blue!50!black,citecolor=blue!50!black,urlcolor=blue!50!black]{hyperref}
\usepackage{caption}
\usepackage{enumitem}

\renewcommand{\arraystretch}{1.15}

\pagestyle{fancy}
\fancyhf{}
\fancyhead[L]{RBF multicu\'adrico con distancia de Gower aplicado a biochar}
\fancyhead[R]{\thepage}
\renewcommand{\headrulewidth}{0.4pt}

\newcommand{\R}{\mathbb{R}}
\newcommand{\Phimat}{\boldsymbol{\Phi}}
\newcommand{\wvec}{\mathbf{w}}
\newcommand{\yvec}{\mathbf{y}}

\title{\Large \textbf{Aplicaci\'on del modelo RBF multicu\'adrico a datos reales
de pir\'olisis catal\'itica de biochar: \\
depuraci\'on del dataset, extensi\'on a variables mixtas mediante \\
distancia de Gower, y resultados para H\textsubscript{2}}}

\author{Informe t\'ecnico}
\date{\today}

\begin{document}

\maketitle

\begin{abstract}
\noindent
Este informe documenta la aplicaci\'on del modelo de aproximaci\'on RBF
multicu\'adrico con selecci\'on \emph{greedy} de centros y b\'usqueda de
hiperpar\'ametros por Latin Hypercube Sampling (desarrollado y validado
previamente sobre datos sint\'eticos) a un conjunto de datos experimentales
reales de pir\'olisis catal\'itica de biochar. Se documentan: (i) el proceso
de depuraci\'on y reorganizaci\'on del dataset original, incluyendo varios
hallazgos relevantes sobre la naturaleza real de las variables (una variable
de entrada no documentada inicialmente, una colinealidad perfecta entre dos
columnas, y la identificaci\'on de tres columnas que result\'aron ser
variables de salida en lugar de entrada); (ii) una extensi\'on metodol\'ogica
necesaria para incorporar dos variables de entrada categ\'oricas (nominales)
mediante una distancia compuesta an\'aloga al coeficiente de similitud de
Gower (Gower, 1971); y (iii) los resultados de aplicar el modelo extendido a
la predicci\'on de la fracci\'on molar de H\textsubscript{2} en el gas de
s\'intesis. Se cierra con una discusi\'on sobre la sensibilidad observada de
los hiperpar\'ametros del modelo frente a la notable estabilidad del
RMSE-LOOCV resultante, conectada con resultados an\'alogos --aunque en un
contexto metodol\'ogicamente distinto-- de la literatura de colocaci\'on RBF
para ecuaciones diferenciales parciales.
\end{abstract}

\tableofcontents

% =====================================================================
\section{Introducci\'on}
% =====================================================================

En un trabajo previo se dise\~n\'o, implement\'o y valid\'o (sobre datos
sint\'eticos con estructura controlada) un modelo de aproximaci\'on mediante
funciones de base radial (RBF) con n\'ucleo multicu\'adrico
$\phi(r)=\sqrt{r^2+c^2}$, incorporando: estandarizaci\'on de variables,
regularizaci\'on de Tikhonov, selecci\'on \emph{greedy} de centros guiada por
validaci\'on cruzada \emph{leave-one-out} (LOOCV) calculada de forma
anal\'itica mediante la f\'ormula de Rippa/Allen, y b\'usqueda de
hiperpar\'ametros mediante Latin Hypercube Sampling (LHS) en lugar de una
grilla regular.

Este informe documenta la primera aplicaci\'on de ese modelo a un conjunto de
datos experimentales reales: un ensayo de pir\'olisis catal\'itica de
biochar, con $N=83$ observaciones. A diferencia del caso sint\'etico
original --en el que las cinco variables independientes eran continuas--
este dataset incluye dos variables de entrada \textbf{categ\'oricas}
(nominales, sin orden): el origen de la biomasa y el tipo de catalizador.
Esto oblig\'o a una extensi\'on metodol\'ogica del modelo, documentada en la
Secci\'on~\ref{sec:gower}.

Antes de llegar a esa extensi\'on, sin embargo, una parte sustancial del
trabajo consisti\'o en depurar y reinterpretar correctamente el dataset
original, cuya estructura conten\'ia varias ambig\"uedades no triviales.
Esa depuraci\'on se documenta en la Secci\'on~\ref{sec:dataset}, porque
varias de sus conclusiones determinan directamente c\'omo se plante\'o el
problema de modelado.

% =====================================================================
\section{El dataset original: estructura y hallazgos del an\'alisis exploratorio}
\label{sec:dataset}
% =====================================================================

\subsection{Estructura del archivo entregado}

El archivo original (\texttt{05\_Raw\_Original\_Dataset\_Reference.xlsx})
conten\'ia 83 filas y 15 columnas, agrupadas por el propio archivo en tres
encabezados: \emph{``ENTRADAS''} (ORIGEN, CAT, CARGA\,\%, TR, N2),
\emph{``OBTENIDAS POR MEDICI\'ON Y TGA''} (PERDIDA AL FUEGO, \% S\'olido,
CENIZA), y \emph{``SALIDAS''} (T, H2, CO, CH4, CO2, H2/CO, PERDIDA+CENIZA).
El usuario indic\'o inicialmente que las variables independientes eran
ORIGEN, CAT, TR, N2, PERDIDA AL FUEGO, \%S\'olido, CENIZA y T (esta \'ultima
mal etiquetada como ``salida'' en el archivo original), y que las variables
dependientes eran H2, CO, CH4 y CO2.

\subsection{Dise\~no experimental subyacente}

El an\'alisis exploratorio revel\'o un dise\~no experimental balanceado:
ORIGEN (amazonas/cimitarra/maceo) y CAT (CaO/MgO/No CAT) forman una grilla
$3\times3=9$ combinaciones, cada una con 9--10 r\'eplicas (Tabla~\ref{tab:diseno}).
Adem\'as, T y TR resultaron tomar \'unicamente 3 valores discretos cada una
(400/600/800\,°C y 0.25/0.5/1.0, respectivamente) por dise\~no experimental,
mientras que N2 y CARGA\,\% var\'ian de forma m\'as continua.

\begin{table}[htbp]
\centering
\caption{Combinaciones ORIGEN $\times$ CAT y n\'umero de r\'eplicas.}
\label{tab:diseno}
\begin{tabular}{@{}llr@{}}
\toprule
ORIGEN & CAT & n \\
\midrule
amazonas & CaO & 9 \\
amazonas & MgO & 9 \\
amazonas & No CAT & 9 \\
cimitarra & CaO & 10 \\
cimitarra & MgO & 9 \\
cimitarra & No CAT & 9 \\
maceo & CaO & 10 \\
maceo & MgO & 9 \\
maceo & No CAT & 9 \\
\bottomrule
\end{tabular}
\end{table}

\subsection{Hallazgo 1: una variable de entrada no documentada}

La columna \texttt{CARGA\,\%} (valores 0, 0.05, 0.10) apareci\'o en el
archivo sin haber sido mencionada como variable independiente. Un detalle
llam\'o la atenci\'on: incluso las filas con \texttt{CAT="No CAT"} (sin
catalizador) presentaban \texttt{CARGA\,\%}$>0$ en algunos casos --es decir,
no se trata simplemente de la fracci\'on de catalizador cargado, que ser\'ia
id\'enticamente cero sin catalizador. Consultado el usuario, se confirm\'o
que \texttt{CARGA\,\%} es, en efecto, una variable de entrada leg\'itima
adicional, independiente de la anterior duda.

\subsection{Hallazgo 2: colinealidad perfecta entre dos columnas}

Se encontr\'o que \texttt{CENIZA} + \texttt{PERDIDA AL FUEGO} $= 1$ de forma
\textbf{exacta} en las 83 filas (correlaci\'on $=-1.000$; diferencia m\'axima
respecto a la columna de referencia \texttt{PERDIDA+CENIZA}, que result\'o
ser id\'enticamente 1, fue cero a precisi\'on de m\'aquina). Ambas columnas
codifican, por tanto, la misma informaci\'on: una es uno menos la otra. Se
decidi\'o retener \'unicamente \texttt{CENIZA} como variable independiente
(en el sentido estad\'istico, no de entrada del modelo -- ver
Hallazgo~4) junto con \texttt{\% S\'olido}, descartando \texttt{PERDIDA AL
FUEGO} por ser enteramente redundante.

\subsection{Hallazgo 3: columnas puramente derivadas}

Se verific\'o que \texttt{H2/CO} $=$ H2 $\div$ CO exactamente (diferencia
m\'axima $=0$) y que \texttt{PERDIDA+CENIZA} $=$ PERDIDA AL FUEGO $+$ CENIZA
exactamente. Ambas son columnas de referencia calculadas a partir de otras ya
presentes en la tabla, sin informaci\'on independiente: no se usan como
entrada ni como salida separada del modelo.

\subsection{Hallazgo 4 (el m\'as relevante): tres columnas son en realidad SALIDAS}
\label{subsec:hallazgo4}

El agrupamiento del propio archivo original ya suger\'ia algo: PERDIDA AL
FUEGO, \% S\'olido y CENIZA no estaban bajo el encabezado ``ENTRADAS'', sino
bajo un encabezado separado, ``OBTENIDAS POR MEDICI\'ON Y TGA''. Un
an\'alisis de correlaci\'on aport\'o evidencia estad\'istica concreta:
\texttt{\% S\'olido} correlaciona con \texttt{T} en $r=-0.60$ (bastante
fuerte). Esa correlaci\'on tiene una explicaci\'on f\'isica inmediata
\textbf{solo} si \texttt{\% S\'olido} es el rendimiento de s\'olido/biochar
de \emph{ese} experimento (a mayor temperatura de pir\'olisis, m\'as
devolatilizaci\'on, menos s\'olido remanente) -- un fen\'omeno bien conocido
en pir\'olisis. Si, en cambio, \texttt{\% S\'olido} fuera una caracterizaci\'on
de la biomasa \emph{cruda} (antes de cualquier tratamiento t\'ermico), no
habr\'ia raz\'on para que correlacionara con la temperatura del proceso.

Puesto en conocimiento de esta evidencia, el usuario confirm\'o: PERDIDA AL
FUEGO, \% S\'olido y CENIZA se miden \textbf{despu\'es} de la pir\'olisis --
caracterizan el s\'olido/biochar resultante, no la biomasa de partida. Son,
por tanto, variables de \textbf{salida}, no de entrada. Usarlas como entrada
para predecir la composici\'on de gas habr\'ia sido una fuga de informaci\'on
(usar el resultado de un experimento para ``predecir'' otro resultado del
mismo experimento), adem\'as de inutilizable en la pr\'actica: esas
caracter\'isticas no se conocen antes de correr el experimento.

Esta correcci\'on cambia sustancialmente el problema de modelado: en lugar de
4 variables de salida (los gases) y 8 de entrada, el problema qued\'o
planteado con \textbf{6 variables de entrada} y \textbf{7 columnas de
salida} (con redundancia entre dos de ellas, ver Hallazgo~2).

\subsection{Hallazgo 5: los cuatro gases son un dato composicional}

Se verific\'o que H2+CO+CH4+CO2 $=1$ exactamente en las 83 filas (desviaci\'on
est\'andar del orden de $10^{-17}$, es decir, precisi\'on de m\'aquina). Los
cuatro gases no son cuatro variables independientes sino una composici\'on
con 3 grados de libertad: conocidos tres de los cuatro valores, el cuarto
queda determinado exactamente. Esto tiene una consecuencia directa sobre la
estrategia de modelado, discutida en la Secci\'on~\ref{subsec:composicional}.

\subsection{Reorganizaci\'on del dataset}

Con base en estos hallazgos, se entreg\'o al usuario un archivo Excel
reorganizado (\texttt{dataset\_biochar\_reorganizado.xlsx}) con las columnas
agrupadas y coloreadas por bloque (entradas; salidas -- gas; salidas --
s\'olido; referencia/derivadas), columnas de verificaci\'on calculadas por
f\'ormula (nunca valores fijos) confirmando ambas restricciones exactas
encontradas, y comentarios document\'ando cada decisi\'on metodol\'ogica
directamente en los encabezados correspondientes.

% =====================================================================
\section{Reformulaci\'on del problema de modelado}
\label{sec:reformulacion}
% =====================================================================

\subsection{Variables de entrada (6)}

\begin{itemize}[nosep]
    \item \textbf{ORIGEN} -- categ\'orica, 3 niveles (amazonas, cimitarra, maceo)
    \item \textbf{CAT} -- categ\'orica, 3 niveles (CaO, MgO, No CAT)
    \item \textbf{CARGA\,\%}, \textbf{TR}, \textbf{N2}, \textbf{T} -- continuas
\end{itemize}

\subsection{Variables de salida y estrategia composicional}
\label{subsec:composicional}

De las 7 columnas de salida disponibles (H2, CO, CH4, CO2, CENIZA, \%
S\'olido, PERDIDA AL FUEGO), se descarta PERDIDA AL FUEGO por redundante
(Hallazgo~2), quedando 6 columnas de las cuales 5 son independientes: 3
grados de libertad en los gases (Hallazgo~5) m\'as CENIZA y \% S\'olido.

Para los gases, se opt\'o por modelar 3 de los 4 directamente v\'ia RBF y
derivar el cuarto por resta ($z_4 = 1-z_1-z_2-z_3$), garantizando as\'i que
cualquier predicci\'on nueva satisfaga exactamente la restricci\'on
composicional --a diferencia de modelar los 4 independientemente y
normalizar al final, o de una transformaci\'on log-ratio (Aitchison), que
se descartaron por mayor complejidad de implementaci\'on e interpretaci\'on
sin una ganancia clara para este caso.

\subsubsection{Qu\'e gas derivar}

Para decidir cu\'al de los cuatro gases conviene derivar (en lugar de
modelar directamente), se realiz\'o un diagn\'ostico r\'apido: regresi\'on
lineal m\'ultiple (entradas continuas estandarizadas + categ\'oricas
\emph{one-hot}) con validaci\'on cruzada de 5 particiones, como proxy
econ\'omico de qu\'e tan predecible es cada gas desde las entradas
(Tabla~\ref{tab:diagnostico_gases}).

\begin{table}[htbp]
\centering
\caption{Diagn\'ostico de predictibilidad lineal (proxy) por gas.}
\label{tab:diagnostico_gases}
\begin{tabular}{@{}lrrr@{}}
\toprule
Gas & $R^2$ CV (lineal) & Rango & Desv. est\'andar \\
\midrule
H2  & 0.925 & 0.009--0.378 & 0.116 \\
CO2 & 0.828 & 0.044--0.551 & 0.131 \\
CH4 & 0.596 & 0.192--0.428 & 0.045 \\
CO  & \textbf{0.367} & 0.100--0.371 & 0.061 \\
\bottomrule
\end{tabular}
\end{table}

CO result\'o, por amplio margen, el m\'as dif\'icil de explicar linealmente
desde las entradas disponibles. Se decidi\'o modelar H2, CH4 y CO2
directamente v\'ia RBF, y derivar CO $=1-$H2$-$CH4$-$CO2. El razonamiento es
doble: (i) al ser CO el m\'as ruidoso/dif\'icil de ajustar directamente,
derivarlo aprovecha la restricci\'on de cierre composicional para
``corregirlo'' con la informaci\'on de los otros tres, en lugar de depender
de un ajuste directo potencialmente inestable; (ii) el radio H2/CO, de
inter\'es pr\'actico seg\'un el archivo original, sigue siendo perfectamente
calculable con el CO derivado. Se advierte, no obstante, que este
diagn\'ostico usa un modelo \textbf{lineal}; una vez el RBF (no lineal) est\'e
ajustado sobre los 4 gases, conviene revisar emp\'iricamente si esta elecci\'on
se mantiene.

Quedan, en definitiva, \textbf{5 modelos RBF independientes} por ajustar:
H2, CH4, CO2, CENIZA y \% S\'olido. Este informe documenta el primero de
ellos (H2) en detalle, como plantilla para los restantes.

% =====================================================================
\section{Extensi\'on metodol\'ogica: distancia mixta tipo Gower}
\label{sec:gower}
% =====================================================================

\subsection{El problema: dos entradas son categ\'oricas}

Todo el modelo RBF descansa en una distancia eucl\'idea
$r=\lVert\mathbf x-\mathbf x'\rVert$ con sentido f\'isico: dos valores de
temperatura cercanos (p.ej. 450\,°C y 500\,°C) deben producir una $r$ chica.
ORIGEN y CAT son variables \textbf{nominales}: no existe ning\'un sentido en
el que ``zeolita'' est\'e m\'as cerca de un catalizador que de otro. Codificar
las categor\'ias como enteros arbitrarios (CaO$=1$, MgO$=2$, No CAT$=3$) le
impondr\'ia al modelo una estructura ordinal inexistente --el modelo
``creer\'ia'' que el catalizador 1 est\'a m\'as cerca del 2 que del 3, sin que
eso signifique nada.

\subsection{El coeficiente de similitud de Gower}

Gower (1971) \cite{gower1971} propuso un coeficiente general de similitud
para conjuntos de datos con variables de tipo mixto (continuas, ordinales,
binarias y nominales), pensado originalmente para taxonom\'ia num\'erica.
Su idea central es promediar, sobre las $p$ variables disponibles, una
disimilitud calculada variable por variable:
\[
d(\mathbf x, \mathbf x') = \frac{1}{p}\sum_{k=1}^{p} d_k(x_k, x_k'),
\]
donde $d_k$ se define seg\'un el tipo de la variable $k$:
\begin{itemize}[nosep]
    \item Si $k$ es continua: $d_k(x_k,x_k') = |x_k-x_k'|/R_k$, con $R_k$ el
    rango observado de la variable (normalizaci\'on a $[0,1]$).
    \item Si $k$ es nominal (nuestro caso): $d_k(x_k,x_k')=0$ si
    $x_k=x_k'$, y $1$ en caso contrario (``\emph{simple matching}'').
\end{itemize}
Gower demostr\'o que la matriz de similitudes resultante es semidefinida
positiva (salvo, posiblemente, en presencia de datos faltantes), lo cual es
relevante porque garantiza que $d$ admite una representaci\'on eucl\'idea
--una propiedad deseable para cualquier distancia que se use dentro de un
n\'ucleo como el multicu\'adrico. Esta construcci\'on es, hoy, el punto de
partida est\'andar en la literatura de an\'alisis de conglomerados y
similitud para datos de tipo mixto \cite{gower1986}.

\subsection{Adaptaci\'on al n\'ucleo multicu\'adrico}

La multicu\'adrica $\phi(r)=\sqrt{r^2+c^2}$ requiere una distancia (o,
equivalentemente, una distancia \emph{al cuadrado} $r^2$, que es la cantidad
que efectivamente se necesita para evaluar $\phi$). Se adapta la idea de
Gower --separar la disimilitud por bloques de variables, seg\'un su tipo--
pero en la versi\'on ``distancia al cuadrado'' (para insertarla directamente
en el n\'ucleo) y reemplazando el promedio con pesos fijos $1/p$ por
\textbf{pesos ajustables} $w_m$, uno por cada variable categ\'orica, tratados
como hiperpar\'ametros del modelo:
\[
\boxed{\;
r^2(\mathbf x,\mathbf x') \;=\; \underbrace{\lVert \tilde{\mathbf x}_{\text{cont}}
- \tilde{\mathbf x}'_{\text{cont}}\rVert^2}_{\text{bloque continuo, z-score}}
\;+\; \sum_{m=1}^{k} w_m \cdot \mathbb{1}\!\left[\,\text{cat}_m \neq \text{cat}_m'\,\right]
\;}
\]
donde $\tilde{\mathbf x}_{\text{cont}}$ es el bloque de variables continuas
ya estandarizado (igual que en el caso puramente continuo), y la suma corre
sobre las $k=2$ variables categ\'oricas (ORIGEN, CAT). Si dos puntos
coinciden en una variable categ\'orica, esa variable no aporta nada a la
distancia; si no coinciden, aporta exactamente $w_m$ --el peso act\'ua como
``cu\'anta distancia (en las mismas unidades que el bloque continuo
estandarizado) representa un cambio de categor\'ia en esa variable''.

A diferencia del coeficiente de Gower original --pensado para medir
similitud descriptivamente, con pesos fijados de antemano por el analista--
aqu\'i los pesos $w_m$ se \textbf{aprenden de los datos}, junto con $c$ y
$\lambda$, minimizando el mismo criterio LOOCV usado para todas las dem\'as
decisiones estructurales del modelo (ver discusi\'on de los dos niveles de
funci\'on objetivo en el informe metodol\'ogico previo).

\subsection{Cola polin\'omica extendida}

La cola polin\'omica (Secci\'on~2.2 del informe metodol\'ogico previo) se
extiende con columnas \emph{dummy} \emph{one-hot} (todas las categor\'ias,
sin omitir ninguna) de cada variable categ\'orica, adem\'as de la constante y
los t\'erminos lineales en las variables continuas. Esto permite capturar el
efecto principal de cada categor\'ia sobre la tendencia global, independiente
de su rol en la distancia radial.

\subsection{Extensi\'on de la b\'usqueda de hiperpar\'ametros}

Con $k=2$ variables categ\'oricas, el espacio de b\'usqueda pasa de dimensi\'on
2 ($c,\lambda$) a dimensi\'on $2+k=4$ ($c,\lambda,w_{\text{ORIGEN}},w_{\text{CAT}}$).
El muestreo mediante Latin Hypercube Sampling (Secci\'on~5 del informe
metodol\'ogico previo) se generaliza de forma directa a esta dimensi\'on
mayor --es, de hecho, precisamente en dimensiones m\'as altas donde LHS
aventaja m\'as claramente a un grid regular (Bergstra \& Bengio, 2012
\cite{bergstra2012}): un grid con la misma resoluci\'on por eje que en el
caso 2D requerir\'ia, en dimensi\'on 4, la cuarta potencia de evaluaciones.

% =====================================================================
\section{Aplicaci\'on a H\textsubscript{2}: refinamiento y resultado final}
\label{sec:h2}
% =====================================================================

\subsection{Primera corrida y hiperpar\'ametros pegados al borde}

Una primera corrida exploratoria (rangos $c\in(0.1,10)$, $\lambda\in(10^{-6},1)$,
$w_m\in(0.01,100)$, 150 puntos LHS) arroj\'o $c=0.156$, $\lambda=2.02\times10^{-6}$
y $w_{\text{ORIGEN}}=0.023$ --los tres sospechosamente cerca de sus respectivos
bordes inferiores (posiciones relativas en escala logar\'itmica de 19\,\%,
5\,\% y 9\,\% del rango, respectivamente). Esto es una se\~nal de alerta
est\'andar en b\'usqueda de hiperpar\'ametros: si el \'optimo encontrado est\'a
en el borde del rango explorado, no se puede descartar que el verdadero
\'optimo est\'e fuera de ese rango.

\subsection{Refinamiento de los rangos}

Se ampliaron los tres rangos hacia los bordes donde se detect\'o el problema,
y se repiti\'o la b\'usqueda con presupuestos crecientes y semillas distintas,
para evaluar tanto la posici\'on del \'optimo como la \textbf{estabilidad}
del resultado entre corridas (Tabla~\ref{tab:refinamiento}).

\begin{table}[htbp]
\centering
\caption{Evoluci\'on de la b\'usqueda de hiperpar\'ametros para H2 a trav\'es
de tres corridas de refinamiento.}
\label{tab:refinamiento}
\small
\begin{tabular}{@{}lrrrrrr@{}}
\toprule
Corrida & Presup. & $c$ & $\lambda$ & $w_{\text{ORIGEN}}$ & $w_{\text{CAT}}$ & RMSE-LOOCV \\
\midrule
1 (rangos iniciales)     & 150 & 0.156  & $2.0\times10^{-6}$ & 0.023 & 9.33 & 0.01712 \\
2 (rangos ampliados)     & 400 & 0.043  & 0.356              & 5.38  & 6.90 & 0.01688 \\
3 (rangos amp., semilla distinta) & 300 & 0.064  & 0.312     & 1.28  & 4.50 & 0.01690 \\
4 (final, mayor presupuesto)      & 500 & 0.081  & 0.324     & 1.28  & 5.00 & 0.01698 \\
\bottomrule
\end{tabular}
\end{table}

Dos observaciones destacan de la Tabla~\ref{tab:refinamiento}. Primero,
$\lambda$ se estabiliz\'o alrededor de $0.3$ en las corridas 2--4 (tras
ampliar su rango hacia arriba), confirmando que su valor inicial cercano al
borde era, efectivamente, un artefacto de un rango demasiado angosto y no un
\'optimo genuino en el l\'imite. Segundo --y m\'as notable--, \textbf{los
hiperpar\'ametros var\'ian sustancialmente entre las corridas 2, 3 y 4 (p.ej.
$w_{\text{ORIGEN}}$ entre 1.28 y 5.38), mientras que el RMSE-LOOCV resultante
permanece casi id\'entico} (0.01688--0.01698, una variaci\'on menor al 1\,\%).
Esta observaci\'on se retoma con detalle en la Secci\'on~\ref{sec:discusion}.

\subsection{Modelo final}

Se adopt\'o como definitivo el resultado de la corrida 4 (mayor presupuesto,
500 puntos LHS): $c=0.0806$, $\lambda=0.324$, $w_{\text{ORIGEN}}=1.28$,
$w_{\text{CAT}}=5.00$, con 29 centros de los 83 candidatos disponibles.
RMSE-LOOCV$=0.0170$; RMSE-entrenamiento$=0.0120$; $R^2$(entrenamiento)$=0.989$.
Dado que $\text{std}(\text{H2})=0.116$, el RMSE-LOOCV implica
$R^2_{\text{LOOCV}}\approx 1-(0.0170/0.116)^2\approx 0.979$: el modelo
explica cerca del 98\,\% de la varianza de H2 en validaci\'on cruzada
\emph{leave-one-out}, sensiblemente mejor que el 92.5\,\% del diagn\'ostico
lineal preliminar (Tabla~\ref{tab:diagnostico_gases}), lo cual sugiere que el
RBF est\'a capturando estructura genuinamente no lineal.

\begin{figure}[htbp]
    \centering
    \includegraphics[width=0.92\textwidth]{figs/diagnosticos_H2.png}
    \caption{Diagn\'osticos del modelo final para H2: convergencia del
    RMSE-LOOCV durante la selecci\'on greedy (arriba-izq.), predicho vs.
    observado (arriba-der.), residuos vs. predicho (abajo-izq.) y
    distribuci\'on de residuos (abajo-der.).}
    \label{fig:diag_h2}
\end{figure}

La Figura~\ref{fig:diag_h2} muestra los diagn\'osticos usuales: convergencia
suave del greedy (deteni\'endose en 29 centros), ajuste muy cercano a la
diagonal en el panel predicho-vs-real, y residuos sin un patr\'on sistem\'atico
fuerte (posiblemente una leve heterocedasticidad, con algo m\'as de dispersi\'on
en el rango medio de predicci\'on).

\subsection{Dependencia parcial e interpretaci\'on fisicoqu\'imica}

\begin{figure}[htbp]
    \centering
    \includegraphics[width=\textwidth]{figs/dependencia_parcial_H2_continuas.png}
    \caption{Efecto de cada variable continua sobre H2 predicho, estratificado
    por tipo de catalizador (ORIGEN fijo en la categor\'ia m\'as frecuente,
    dem\'as variables continuas en su media muestral).}
    \label{fig:pd_cont}
\end{figure}

\begin{figure}[htbp]
    \centering
    \includegraphics[width=0.75\textwidth]{figs/dependencia_parcial_H2_categoricas.png}
    \caption{Efecto principal de ORIGEN y CAT sobre H2 predicho (resto de
    variables en su valor de referencia/media).}
    \label{fig:pd_cat}
\end{figure}

Dado que $w_{\text{CAT}}$ domina consistentemente sobre $w_{\text{ORIGEN}}$
en todas las corridas de la Tabla~\ref{tab:refinamiento}, las curvas de
dependencia parcial de las variables continuas se estratificaron por CAT
(Figura~\ref{fig:pd_cont}). Los resultados son f\'isicoqu\'imicamente
coherentes:

\begin{itemize}
    \item \textbf{La temperatura (T) domina el comportamiento}: H2 predicho
    sube de $\sim$0.10 a 400\,°C hasta $\sim$0.35--0.42 a 800\,°C, con
    diferencia el efecto m\'as fuerte de las seis entradas. Es consistente
    con el craqueo secundario y las reacciones de reformado que se
    intensifican a mayor temperatura de pir\'olisis.
    \item \textbf{CaO $>$ MgO $>$ No CAT}, de forma consistente en todo el
    rango de T. Coincide con la qu\'imica esperada: el CaO es conocido por
    promover reacciones de \emph{water-gas-shift} y reformado
    intensificado por sorci\'on (captura de CO2 como CaCO$_3$, desplazando
    el equilibrio hacia m\'as H2), un mecanismo bien documentado en
    pir\'olisis catal\'itica de biomasa.
    \item CARGA\,\%, TR y N2 muestran efectos mucho m\'as discretos que T.
    \item El efecto principal de ORIGEN result\'o de magnitud comparable al
    de CAT en el gr\'afico de barras (Figura~\ref{fig:pd_cat}, rango
    $\sim$0.22--0.28 para ambas variables) --pese a que los pesos
    $w_{\text{ORIGEN}}$ y $w_{\text{CAT}}$ tomaron valores bastante distintos
    entre corridas (Tabla~\ref{tab:refinamiento}). Esto refuerza que los
    pesos categ\'oricos individuales no son muy identificables con $N=83$
    (superficie plana en esa direcci\'on del espacio de hiperpar\'ametros),
    pero las \emph{predicciones} del modelo final son consistentes.
\end{itemize}

Que el modelo recupere, sin haber sido guiado expl\'icitamente hacia ello,
un ordenamiento CaO $>$ MgO $>$ No CAT consistente con la literatura de
pir\'olisis catal\'itica es la mejor validaci\'on cualitativa disponible en
esta etapa: sugiere que el RBF est\'a aprendiendo estructura f\'isica real,
no ruido.

% =====================================================================
\section{Discusi\'on: sensibilidad de hiperpar\'ametros vs.\ estabilidad del LOOCV}
\label{sec:discusion}
% =====================================================================

\subsection{La observaci\'on emp\'irica}

La Tabla~\ref{tab:refinamiento} documenta un patr\'on que merece discusi\'on
expl\'icita: entre las corridas 2, 3 y 4, los hiperpar\'ametros individuales
--particularmente $c$ (var\'ia por un factor $\sim$2) y $w_{\text{ORIGEN}}$
(var\'ia por un factor $>4$)-- cambian de forma sustancial seg\'un la
semilla y el presupuesto de b\'usqueda usados, mientras que el
RMSE-LOOCV resultante permanece pr\'acticamente constante (dentro de un
0.6\,\% de variaci\'on). Esto sugiere que la superficie de RMSE-LOOCV en
funci\'on de $(c,\lambda,w_{\text{ORIGEN}},w_{\text{CAT}})$ es, en una
regi\'on amplia del espacio de b\'usqueda, relativamente \textbf{plana}:
muchas combinaciones distintas de hiperpar\'ametros producen un desempe\~no
predictivo casi id\'entico.

\subsection{Conexi\'on con la literatura de colocaci\'on RBF para EDPs}

Esta observaci\'on resuena con una discusi\'on de larga data en la
literatura de m\'etodos RBF de colocaci\'on para ecuaciones diferenciales
parciales (EDPs) --un uso de las funciones de base radial metodol\'ogicamente
distinto al de este trabajo (colocaci\'on/interpolaci\'on de una EDP, en
lugar de aproximaci\'on regularizada de datos ruidosos), pero que comparte
el mismo n\'ucleo matem\'atico y, por tanto, los mismos hiperpar\'ametros de
forma. Kansa (1990) \cite{kansa1990a,kansa1990b} introdujo el uso de la
multicu\'adrica para resolver EDPs por colocaci\'on directa (el llamado
\emph{m\'etodo de Kansa}), notando desde entonces que el par\'ametro de
forma $c$ afecta fuertemente tanto la precisi\'on como el n\'umero de
condici\'on del sistema lineal resultante.

Schaback (1995) \cite{schaback1995} formaliz\'o esta tensi\'on como un
\emph{``principio de incertidumbre''}: en el m\'etodo directo de resoluci\'on
del sistema RBF (RBF-Direct), no es posible mejorar la precisi\'on
(t\'ipicamente lograda con par\'ametros de forma m\'as chicos, es decir,
n\'ucleos m\'as ``planos'') sin empeorar simult\'aneamente el n\'umero de
condici\'on de la matriz de interpolaci\'on. Sin embargo, trabajos
posteriores (Fornberg y Wright, y Wright y Fornberg, 2017
\cite{wright2017}) mostraron que este mal condicionamiento es un
\textbf{artefacto del algoritmo de resoluci\'on directo}, no una limitaci\'on
fundamental de la interpolaci\'on RBF en s\'i: m\'etodos alternativos y
num\'ericamente estables (RBF-QR, Contour-Pad\'e, RBF-GA) permiten evaluar
el interpolante para \emph{cualquier} valor del par\'ametro de forma,
incluido el l\'imite ``plano'' ($c\to\infty$ en la parametrizaci\'on de este
informe), sin la explosi\'on del n\'umero de condici\'on que sufre el
m\'etodo directo. Kansa y Carlson (1992) \cite{kansa1992} exploraron adem\'as
el uso de par\'ametros de forma \textbf{variables} (distintos por centro), en
lugar de un \'unico valor global, encontrando mejoras de precisi\'on
adicionales.

\subsection{El comentario del usuario: experiencia previa en colocaci\'on RBF para EDPs}

El usuario aport\'o, a partir de su propia experiencia previa aplicando
colocaci\'on RBF a la soluci\'on de EDPs, una observaci\'on directamente
relevante: en ese contexto (interpolaci\'on/colocaci\'on, no aproximaci\'on de
datos ruidosos), la literatura muestra que los hiperpar\'ametros del n\'ucleo
($c$, y por extensi\'on otros par\'ametros de escala) no son especialmente
determinantes para la precisi\'on final de la soluci\'on, siempre que se
manejen dentro de rangos razonablemente amplios; la precisi\'on de la
soluci\'on est\'a gobernada, principalmente, por el \textbf{n\'umero de
centros (nodos)} usados para discretizar el dominio.

Esta observaci\'on es coherente con un resultado bien establecido en la
teor\'ia de aproximaci\'on de funciones mediante RBF: la convergencia del
interpolante, a medida que aumenta la densidad de centros (equivalentemente,
a medida que la distancia de llenado o \emph{fill distance} $h\to0$), est\'a
gobernada principalmente por la suavidad del n\'ucleo y por $h$, mientras que
el par\'ametro de forma afecta la \textbf{constante} y el \textbf{r\'egimen
num\'erico} (condicionamiento) de esa convergencia, pero no su existencia
\cite{wendland2005,fasshauer2007}. En otras palabras: dado un n\'umero
suficiente de nodos bien distribuidos, existe t\'ipicamente un rango amplio
de valores de $c$ que producen una precisi\'on comparable y cercana al
\'optimo; el verdadero cuello de botella para la precisi\'on es cu\'antos
nodos se usan y c\'omo est\'an distribuidos, no el valor preciso de $c$
dentro de ese rango razonable.

\subsection{S\'intesis y diferencias importantes a se\~nalar}

La observaci\'on emp\'irica de la Secci\'on~\ref{sec:discusion} (variaci\'on
sustancial de hiperpar\'ametros con RMSE-LOOCV casi constante) es, en
esp\'iritu, consistente con esta literatura: en ambos casos, una vez que el
n\'ucleo opera en un r\'egimen num\'ericamente razonable, el resultado final
es relativamente insensible al valor preciso de los hiperpar\'ametros de
forma, y lo que realmente gobierna la calidad del resultado es la asignaci\'on
de centros/nodos.

Conviene, sin embargo, se\~nalar dos diferencias metodol\'ogicas que
limitan la extensi\'on directa de esta analog\'ia:

\begin{enumerate}
    \item \textbf{Interpolaci\'on/colocaci\'on vs.\ aproximaci\'on
    regularizada de datos ruidosos.} En colocaci\'on de EDPs, $N$ (el n\'umero
    de nodos) es una elecci\'on de dise\~no del usuario, libremente
    incrementable como parte de un refinamiento de malla, y la convergencia
    se estudia expl\'icitamente como funci\'on de $N\to\infty$. En este
    trabajo, $N=83$ es fijo (proviene de un ensayo experimental, no se puede
    ``agregar m\'as nodos'' a voluntad), y el n\'umero de centros efectivamente
    usados por el modelo no es una elecci\'on libre del usuario sino el
    resultado del criterio de parada temprana del Algoritmo greedy guiado
    por LOOCV (Secci\'on~2.8 del informe metodol\'ogico previo). La
    ``regulaci\'on por n\'umero de centros'' que menciona la literatura de
    EDPs ocurre aqu\'i de forma \textbf{autom\'atica} v\'ia ese criterio, no
    como un par\'ametro de refinamiento que se ajuste directamente.

    \item \textbf{Ruido.} La colocaci\'on de EDPs t\'ipicamente opera sobre
    datos exactos (la propia EDP y sus condiciones de frontera); aqu\'i los
    datos son observaciones experimentales ruidosas, y la regularizaci\'on
    $\lambda$ juega un rol adicional (controlar sobreajuste al ruido) que no
    tiene un an\'alogo directo en colocaci\'on pura.
\end{enumerate}

A pesar de estas diferencias, la convergencia cualitativa de ambas l\'ineas
de evidencia --la de este trabajo y la de la literatura de EDPs citada por
el usuario-- resulta \'util como una segunda validaci\'on, independiente del
diagn\'ostico fisicoqu\'imico de la Secci\'on~\ref{sec:h2}, de que el
comportamiento observado (baja sensibilidad de los hiperpar\'ametros de
forma, en un r\'egimen num\'ericamente razonable) no es un artefacto
espec\'ifico de este dataset, sino una propiedad m\'as general de los
m\'etodos RBF.

% =====================================================================
\section{Limitaciones y pr\'oximos pasos}
% =====================================================================

\begin{itemize}
    \item El diagn\'ostico de qu\'e gas derivar (Secci\'on~\ref{subsec:composicional})
    us\'o un modelo lineal; conviene revisarlo una vez ajustados los 3
    modelos RBF de gases, comparando el error real de CO derivado vs.\ CO
    modelado directamente.
    \item Falta replicar el procedimiento completo (b\'usqueda de rangos +
    ajuste final + diagn\'osticos) para CH4, CO2, CENIZA y \% S\'olido,
    usando la configuraci\'on de la Secci\'on~\ref{sec:h2} como plantilla.
    \item La posible correlaci\'on moderada entre TR y N2 ($r=0.81$,
    reportada en el informe metodol\'ogico previo) no se ha investigado en
    profundidad en el contexto del dataset real; conviene revisarla al
    completar los 5 modelos.
    \item Como se discuti\'o en la Secci\'on~\ref{sec:discusion}, los pesos
    categ\'oricos individuales no parecen bien identificados con $N=83$; no
    se recomienda interpretarlos como una medida precisa de ``importancia
    relativa'' de ORIGEN vs.\ CAT, m\'as all\'a de la conclusi\'on cualitativa
    de que ambos importan de forma no despreciable.
\end{itemize}

% =====================================================================
\section{Conclusiones}
% =====================================================================

Se aplic\'o exitosamente el modelo RBF multicu\'adrico con selecci\'on
greedy y b\'usqueda de hiperpar\'ametros por LHS a un dataset experimental
real de pir\'olisis catal\'itica de biochar, tras un proceso de depuraci\'on
que revel\'o correcciones sustanciales al planteamiento inicial del problema
(una variable de entrada no documentada, una colinealidad perfecta, y --lo
m\'as importante-- la identificaci\'on de tres columnas que en realidad son
salidas del proceso, no entradas). Se extendi\'o el modelo para incorporar
variables categ\'oricas mediante una distancia compuesta an\'aloga al
coeficiente de similitud de Gower (1971), con pesos categ\'oricos ajustados
como hiperpar\'ametros adicionales v\'ia LHS.

El modelo resultante para H2 alcanza un RMSE-LOOCV de 0.0170
($R^2_{\text{LOOCV}}\approx0.98$), y --sin haber sido guiado expl\'icitamente
hacia ello-- recupera un ordenamiento de efecto catal\'itico (CaO $>$ MgO
$>$ No CAT) consistente con la qu\'imica de pir\'olisis catal\'itica conocida,
adem\'as de un fuerte efecto de la temperatura consistente con el craqueo
secundario esperado. La notable estabilidad del RMSE-LOOCV frente a
variaciones sustanciales en los hiperpar\'ametros individuales, observada
emp\'iricamente en este trabajo, resuena con hallazgos an\'alogos --aunque en
un contexto metodol\'ogico distinto-- en la literatura de colocaci\'on RBF
para ecuaciones diferenciales parciales, donde se ha documentado que, dentro
de un r\'egimen num\'ericamente razonable, la precisi\'on final est\'a
gobernada principalmente por el n\'umero y distribuci\'on de los centros, no
por el valor preciso de los par\'ametros de forma.

% =====================================================================
\begin{thebibliography}{99}

\bibitem{gower1971}
Gower, J.\ C.\ (1971).
\emph{A general coefficient of similarity and some of its properties}.
Biometrics, 27(4), 857--871.

\bibitem{gower1986}
Gower, J.\ C., \& Legendre, P.\ (1986).
\emph{Metric and Euclidean properties of dissimilarity coefficients}.
Journal of Classification, 3(1), 5--48.

\bibitem{bergstra2012}
Bergstra, J., \& Bengio, Y.\ (2012).
\emph{Random search for hyper-parameter optimization}.
Journal of Machine Learning Research, 13, 281--305.

\bibitem{kansa1990a}
Kansa, E.\ J.\ (1990).
\emph{Multiquadrics---A scattered data approximation scheme with
applications to computational fluid dynamics---I: surface approximations
and partial derivative estimates}.
Computers \& Mathematics with Applications, 19(8--9), 127--145.

\bibitem{kansa1990b}
Kansa, E.\ J.\ (1990).
\emph{Multiquadrics---A scattered data approximation scheme with
applications to computational fluid dynamics---II: solutions to parabolic,
hyperbolic and elliptic partial differential equations}.
Computers \& Mathematics with Applications, 19(8--9), 147--161.

\bibitem{kansa1992}
Kansa, E.\ J., \& Carlson, R.\ E.\ (1992).
\emph{Improved accuracy of multiquadric interpolation using variable shape
parameters}.
Computers \& Mathematics with Applications, 24(12), 99--120.

\bibitem{schaback1995}
Schaback, R.\ (1995).
\emph{Error estimates and condition numbers for radial basis function
interpolation}.
Advances in Computational Mathematics, 3(3), 251--264.

\bibitem{wright2017}
Wright, G.\ B., \& Fornberg, B.\ (2017).
\emph{Stable computations with flat radial basis functions using
vector-valued rational approximations}.
Journal of Computational Physics, 331, 137--156.

\bibitem{wendland2005}
Wendland, H.\ (2005).
\emph{Scattered Data Approximation}.
Cambridge University Press.

\bibitem{fasshauer2007}
Fasshauer, G.\ E.\ (2007).
\emph{Meshfree Approximation Methods with MATLAB}.
World Scientific.

\end{thebibliography}

\end{document}
